\chapter{The 2D Phase-Plane Framework: Validation}
\label{ch:phase-plane}

\Cref{ch:tl-horizontal,ch:tl-vertical} showed that simple amplitude
differencing of migrated images detects sub-wavelength displacement only down
to a finite floor. This chapter applies the phase-plane shift estimator
derived in \cref{sec:th-fourier-shift,sec:th-wls,sec:th-material-change} and
implemented as described in \cref{sec:meth-phaseplane} to the same horizontal
and vertical displacement datasets, plus a dedicated noise-robustness test,
and shows that it remains accurate at displacement scales where the amplitude
image is already featureless. It then complements the global wavenumber-plane
fit with the local, trace-based views of \cref{sec:th-local} ---
instantaneous phase imaging, spectral-line fitting, cross-phase
spectrograms, and localised-STFT phase decomposition --- to show that all of
these independent analyses agree.

Because several of the analyses below are repeated identically across seven
(or six) wavelength scales, only one or two representative scales are shown
in this chapter; the complete sweep for every analysis is given in
\cref{app:extended-sweeps}, so that every one of the 81 figures produced by
\texttt{TimeLapse\_Processing.ipynb} for these scenarios appears at least
once in the thesis (the three fluid-front phase-plane figures from the same
notebook, TLP\_013--015, are presented in \cref{ch:fluidflow} instead, where
they belong thematically).

\section{Migrated Baseline Overview Across Scenarios}
\label{sec:tlp-overview}

\Cref{fig:tlp-overview} shows the starting point for every analysis in this
chapter: the raw background-subtracted B-scans and the three migrated images
(Kirchhoff, Gazdag, back-propagation) for every scenario used below.

\begin{figure}[htbp]
  \centering
  \begin{subfigure}[t]{0.48\textwidth}
    \includegraphics[width=\linewidth]{TLP_001_Raw_Background-Subtracted_B-Scans__All_Scenarios.png}
    \caption{Raw, background-subtracted}
  \end{subfigure}
  \hfill
  \begin{subfigure}[t]{0.48\textwidth}
    \includegraphics[width=\linewidth]{TLP_002_Kirchhoff_Migration__All_Scenarios.png}
    \caption{Kirchhoff-migrated}
  \end{subfigure}

  \vspace{0.6em}
  \begin{subfigure}[t]{0.48\textwidth}
    \includegraphics[width=\linewidth]{TLP_003_Gazdag_Migration__All_Scenarios.png}
    \caption{Gazdag-migrated}
  \end{subfigure}
  \hfill
  \begin{subfigure}[t]{0.48\textwidth}
    \includegraphics[width=\linewidth]{TLP_004_Back-prop_Migration__All_Scenarios.png}
    \caption{Back-propagation-migrated}
  \end{subfigure}
  \caption{Overview of all scenarios re-used from \cref{ch:tl-horizontal} in
    this chapter: (a) raw background-subtracted B-scans; (b--d) the same
    data migrated with Kirchhoff, Gazdag, and back-propagation migration.}
  \label{fig:tlp-overview}
\end{figure}

\section{Phase-Plane Fit: Horizontal Shift Validation}
\label{sec:tlp-horizontal-validation}

For every lateral-displacement scenario of \cref{ch:tl-horizontal}, the
estimator of \cref{sec:meth-phaseplane} is applied to the baseline/monitor
pair, cropped to $\pm2.5\lambda$ around the scatterer apex (located from the
Hilbert envelope of the baseline image). \Cref{fig:tlp-horiz-validation}
shows the recovered $(\Dz,\Dx)$ compared with the known ground truth for all
three migration methods.

\begin{figure}[htbp]
  \centering
  \begin{subfigure}[t]{0.32\textwidth}
    \includegraphics[width=\linewidth]{TLP_005_Phase-plane_shift_estimation__Kirchhoff____Baseline_vs_each.png}
    \caption{Kirchhoff}
  \end{subfigure}
  \hfill
  \begin{subfigure}[t]{0.32\textwidth}
    \includegraphics[width=\linewidth]{TLP_006_Phase-plane_shift_estimation__Gazdag____Baseline_vs_each_sce.png}
    \caption{Gazdag}
  \end{subfigure}
  \hfill
  \begin{subfigure}[t]{0.32\textwidth}
    \includegraphics[width=\linewidth]{TLP_007_Phase-plane_shift_estimation__Back-prop____Baseline_vs_each.png}
    \caption{Back-propagation}
  \end{subfigure}
  \caption{2D phase-plane shift estimation applied to the baseline-versus-each-scenario
    pairs of \cref{ch:tl-horizontal}, for all three migration methods.}
  \label{fig:tlp-horiz-validation}
\end{figure}

\draftnote{read off \cref{fig:tlp-horiz-validation} the smallest lateral
displacement at which the estimated $\Dx$ still tracks the true value to
within (state your accepted tolerance, e.g.\ $\pm 5\,\%$ or $\pm\SI{1}{mm}$),
for each migration method, and contrast this explicitly with the amplitude
floor found in \cref{sec:tl-detectability}.}

\section{Robustness Under Noise}
\label{sec:tlp-noise}

The same estimator is applied to the noisy lateral-displacement dataset of
\cref{sec:tl-noise} (\cref{fig:tlp-noise}), for the two migration methods
available on the noisy data (Kirchhoff and Gazdag).

\begin{figure}[htbp]
  \centering
  \begin{subfigure}[t]{0.48\textwidth}
    \includegraphics[width=\linewidth]{TLP_008_NoisyTimeLapse__Phase-plane_shift_estimation__Kirchhoff____B.png}
    \caption{Kirchhoff}
  \end{subfigure}
  \hfill
  \begin{subfigure}[t]{0.48\textwidth}
    \includegraphics[width=\linewidth]{TLP_009_NoisyTimeLapse__Phase-plane_shift_estimation__Gazdag____Base.png}
    \caption{Gazdag}
  \end{subfigure}
  \caption{2D phase-plane shift estimation applied to the \emph{noisy}
    lateral-displacement dataset of \cref{sec:tl-noise}.}
  \label{fig:tlp-noise}
\end{figure}

\draftnote{quantify, for both methods, how much the additive Laplace noise
degrades the estimated $\Dx$ relative to the clean result of
\cref{fig:tlp-horiz-validation}, and relate this to the amplitude-weighting
argument of \cref{sec:th-mask-weight} --- this figure pair is the only
direct evidence currently available for the noise-resilience claim
discussed qualitatively in \cref{ch:discussion}.}

\section{Phase-Plane Fit: Vertical Shift Validation}
\label{sec:tlp-vertical}

\Cref{fig:tlp-vert-validation} repeats \cref{sec:tlp-horizontal-validation}
for the vertical-displacement dataset of \cref{ch:tl-vertical}, with the ROI
crop and cross-spectrum fit unchanged except that the displacement being
recovered is now $\Dz$ rather than $\Dx$.

\begin{figure}[htbp]
  \centering
  \begin{subfigure}[t]{0.32\textwidth}
    \includegraphics[width=\linewidth]{TLP_010_VerticalTimeLapse__Phase-plane_shift_estimation__Kirchhoff.png}
    \caption{Kirchhoff}
  \end{subfigure}
  \hfill
  \begin{subfigure}[t]{0.32\textwidth}
    \includegraphics[width=\linewidth]{TLP_011_VerticalTimeLapse__Phase-plane_shift_estimation__Gazdag____B.png}
    \caption{Gazdag}
  \end{subfigure}
  \hfill
  \begin{subfigure}[t]{0.32\textwidth}
    \includegraphics[width=\linewidth]{TLP_012_VerticalTimeLapse__Phase-plane_shift_estimation__Back-prop.png}
    \caption{Back-propagation}
  \end{subfigure}
  \caption{2D phase-plane shift estimation applied to the baseline-versus-each-scenario
    pairs of \cref{ch:tl-vertical}, for all three migration methods.}
  \label{fig:tlp-vert-validation}
\end{figure}

\draftnote{state whether the vertical phase-plane fit remains accurate to
the same smallest displacement found for the lateral case in
\cref{sec:tlp-horizontal-validation}, or whether the lateral/vertical
asymmetry predicted in \cref{sec:th-duality} (a constant vertical
\emph{plateau} versus a lateral \emph{slope}) shows up as a difference in
achievable precision between \cref{fig:tlp-horiz-validation,fig:tlp-vert-validation}.}

\section{Instantaneous Phase Imaging --- Lateral Displacement}
\label{sec:tlp-instphase-horizontal}

As a local, spatial-domain complement to the global plane fit above, the 2D
Riesz transform of \cref{eq:kx-inst}--\cref{eq:dphi-lateral} is applied
directly to the Gazdag-migrated baseline and monitor images, producing an
instantaneous-phase difference map $\Dphi(z,x)$ and, at the scatterer depth,
a 1D phase cross-section whose slope at the zero-crossing encodes $\Dx$. This
analysis is repeated at all seven lateral scales from $2\lambda$ to
$\tfrac{1}{32}\lambda$; \cref{fig:tlp-instphase-horiz-example} shows the
representative $\tfrac14\lambda$ case ($\Dx=\SI{280}{mm}$), and
\cref{fig:tlp-instphase-horiz-summary} summarises the fitted zero-crossing
slope across all seven scales. The full seven-scale sweep is given in
\cref{app:sec-instphase-horizontal}.

\begin{figure}[htbp]
  \centering
  \begin{subfigure}[t]{0.32\textwidth}
    \includegraphics[width=\linewidth]{TLP_025_Instantaneous_Phase__Gazdag____Baseline_vs_¼λ___Δx__280_mm.png}
    \caption{Instantaneous phase}
  \end{subfigure}
  \hfill
  \begin{subfigure}[t]{0.32\textwidth}
    \includegraphics[width=\linewidth]{TLP_026_Phase_cross-section__z__676_mm____Gazdag____¼λ__Δx__280_mm.png}
    \caption{Cross-section, $z=\SI{676}{mm}$}
  \end{subfigure}
  \hfill
  \begin{subfigure}[t]{0.32\textwidth}
    \includegraphics[width=\linewidth]{TLP_027_Phase_cross-section__zoomed__½λ_around_scatterers.png}
    \caption{Cross-section, zoomed}
  \end{subfigure}
  \caption{Instantaneous-phase analysis for the representative
    $\Dx=\tfrac14\lambda=\SI{280}{mm}$ lateral displacement, Gazdag-migrated:
    (a) the 2D wrapped phase-difference map; (b) the 1D phase cross-section
    at the scatterer depth $z=\SI{676}{mm}$; (c) the same cross-section
    zoomed to $\pm\tfrac12\lambda$ around the scatterer.}
  \label{fig:tlp-instphase-horiz-example}
\end{figure}

\begin{figure}[htbp]
  \centering
  \includegraphics[width=0.7\linewidth]{TLP_037_Gazdag__Δφ_zero-crossing_slope_vs_scatterer_separation.png}
  \caption{Fitted zero-crossing slope of the lateral instantaneous
    phase-difference cross-section, as a function of true scatterer
    displacement, across all seven scales from $2\lambda$ to
    $\tfrac{1}{32}\lambda$ (Gazdag-migrated).}
  \label{fig:tlp-instphase-horiz-summary}
\end{figure}

\section{Instantaneous Phase Imaging --- Vertical Displacement}
\label{sec:tlp-instphase-vertical}

The same instantaneous-phase analysis is repeated for the vertical
displacement dataset, with the cross-section now taken along $z$ at the
fixed lateral position $x=\SI{2.0}{m}$ of the scatterer, across six scales
from $1\lambda$ to $\tfrac{1}{32}\lambda$. \Cref{fig:tlp-instphase-vert-example}
shows the representative $\tfrac14\lambda$ case ($\Dz=\SI{280}{mm}$), and
\cref{fig:tlp-instphase-vert-summary} summarises the mean phase difference
across all six scales. The full sweep is given in
\cref{app:sec-instphase-vertical}.

\begin{figure}[htbp]
  \centering
  \begin{subfigure}[t]{0.32\textwidth}
    \includegraphics[width=\linewidth]{TLP_044_VerticalTimeLapse__Instantaneous_Phase__Gazdag____Baseline_v.png}
    \caption{Instantaneous phase}
  \end{subfigure}
  \hfill
  \begin{subfigure}[t]{0.32\textwidth}
    \includegraphics[width=\linewidth]{TLP_045_Phase_cross-section__x__2000_mm____Gazdag____¼λ__Δz__280_mm.png}
    \caption{Cross-section, $x=\SI{2000}{mm}$}
  \end{subfigure}
  \hfill
  \begin{subfigure}[t]{0.32\textwidth}
    \includegraphics[width=\linewidth]{TLP_046_Phase_cross-section__zoomed__½λ_around_scatterer_depths.png}
    \caption{Cross-section, zoomed}
  \end{subfigure}
  \caption{Instantaneous-phase analysis for the representative
    $\Dz=\tfrac14\lambda=\SI{280}{mm}$ vertical displacement, Gazdag-migrated:
    (a) the 2D wrapped phase-difference map; (b) the 1D phase cross-section
    at $x=\SI{2000}{mm}$; (c) the same cross-section zoomed to
    $\pm\tfrac12\lambda$ around the scatterer depths.}
  \label{fig:tlp-instphase-vert-example}
\end{figure}

\begin{figure}[htbp]
  \centering
  \includegraphics[width=0.7\linewidth]{TLP_056_VerticalTimeLapse__Gazdag__Mean_Δφ_vs_scatterer_vertical_shi.png}
  \caption{Mean instantaneous phase difference $\Dphi$ as a function of true
    vertical scatterer shift, across all six scales from $1\lambda$ to
    $\tfrac{1}{32}\lambda$ (Gazdag-migrated). Unlike the lateral case
    (\cref{fig:tlp-instphase-horiz-summary}), this is a plateau level rather
    than a cross-section slope, consistent with the asymmetry derived in
    \cref{sec:th-duality}.}
  \label{fig:tlp-instphase-vert-summary}
\end{figure}

\section{Spectral-Line Phase Analysis}
\label{sec:tlp-spectral-line}

The localised-Fourier-shift spectral-line equation
\cref{eq:local-phase-line} is verified directly on individual traces using
the CLSSA decomposition of \texttt{phase\_decomposition.py}, applied both to
the Gazdag-migrated image and to the raw, unmigrated B-scan, at the column
position containing each scenario's scatterer. \Cref{fig:tlp-spectral-line-example}
compares the representative $\tfrac14\lambda$ case in both domains; the full
seven-scale sweep for both is given in
\cref{app:sec-spectral-line}.

\begin{figure}[htbp]
  \centering
  \begin{subfigure}[t]{0.48\textwidth}
    \includegraphics[width=\linewidth]{TLP_060_Spectral_Line_CLSSA____Gazdag____¼λ___Δx__280_mm__02500λ.png}
    \caption{Gazdag-migrated}
  \end{subfigure}
  \hfill
  \begin{subfigure}[t]{0.48\textwidth}
    \includegraphics[width=\linewidth]{TLP_067_Spectral_Line_CLSSA____Raw_Unmigrated____¼λ___Δx__280_mm__02.png}
    \caption{Raw, unmigrated}
  \end{subfigure}
  \caption{CLSSA spectral-line decomposition (amplitude spectrum, phase
    spectrum, and phase gather, \cref{sec:th-local}) at the representative
    $\Dx=\tfrac14\lambda=\SI{280}{mm}$ scale: (a) Gazdag-migrated trace; (b)
    raw, unmigrated trace at the same scenario.}
  \label{fig:tlp-spectral-line-example}
\end{figure}

\draftnote{state whether the migrated and raw spectral lines in
\cref{fig:tlp-spectral-line-example} give consistent $\Dt$/slope estimates
once converted with \cref{eq:dt-dz}, which would confirm that migration is
not required for the temporal spectral-line method to work, only for the
2D wavenumber plane fit of \cref{sec:th-wls}.}

\section{Cross-Phase Spectrograms}
\label{sec:tlp-spectrogram}

\Cref{fig:tlp-spectrogram-example} shows the cross-phase spectrogram
$\Dphi(\tau,f)$ of \cref{sec:th-local} for two contrasting scales: a large,
easily visible $\tfrac14\lambda$ displacement and the smallest,
$\tfrac{1}{32}\lambda$, where the coherent window around the scatterer's
two-way time becomes much harder to distinguish from the incoherent
background. The full seven-scale sweep is given in
\cref{app:sec-spectrogram}.

\begin{figure}[htbp]
  \centering
  \begin{subfigure}[t]{0.48\textwidth}
    \includegraphics[width=\linewidth]{TLP_074_Cross-Phase_Spectrogram_CLSSA__ΔΦτf____Gazdag____¼λ.png}
    \caption{$\Dx=\tfrac14\lambda$}
  \end{subfigure}
  \hfill
  \begin{subfigure}[t]{0.48\textwidth}
    \includegraphics[width=\linewidth]{TLP_077_Cross-Phase_Spectrogram_CLSSA__ΔΦτf____Gazdag____¹₃₂λ.png}
    \caption{$\Dx=\tfrac{1}{32}\lambda$}
  \end{subfigure}
  \caption{Cross-phase spectrogram $\Dphi(\tau,f)$ (two-way time against
    frequency) at two contrasting lateral displacement scales,
    Gazdag-migrated.}
  \label{fig:tlp-spectrogram-example}
\end{figure}

\section{Localised Fourier Shift via Short-Time Fourier Transform}
\label{sec:tlp-stft}

Finally, the three-panel localised-STFT decomposition of
\cref{eq:local-shift,eq:local-phase-line} --- amplitude spectrum
$A(\tau,f)$, phase spectrum $\Dphi(\tau,f)$, and phase-angle domain
$A(\tau,\theta)$ --- is applied with a Gaussian analysis window.
\Cref{fig:tlp-stft-example} shows two of the seven outputs produced by this
sweep; the full set is given in \cref{app:sec-stft}.

\begin{figure}[htbp]
  \centering
  \begin{subfigure}[t]{0.48\textwidth}
    \includegraphics[width=\linewidth]{TLP_078_Localized_Fourier_Shift_STFT_Gaussian_win12_smp____Gazdag.png}
    \caption{Scale 1 of 7}
  \end{subfigure}
  \hfill
  \begin{subfigure}[t]{0.48\textwidth}
    \includegraphics[width=\linewidth]{TLP_081_Localized_Fourier_Shift_STFT_Gaussian_win12_smp____Gazdag.png}
    \caption{Scale 4 of 7}
  \end{subfigure}
  \caption{Localised-STFT phase decomposition (Gaussian window, Gazdag-migrated)
    for two of the seven lateral displacement scales.
    \draftnote{TLP\_078--084 all share an identical auto-generated filename
    stem and do not encode which $\Dx$ each panel corresponds to --- confirm
    the scale for each of the seven figures against the
    \texttt{TimeLapse\_Processing.ipynb} cell order before finalising these
    captions.}}
  \label{fig:tlp-stft-example}
\end{figure}
